\section{From a valid record to a financial transition}
\label{sec:typed-return-transitions}

A company can have a balanced opening book and enough information to assess
the conditions of a method without yet having a journal for the next event.
The method must determine the amounts, identify the accounts, and change the
contract and inventory records consistently. The return example makes these
steps concrete. At the settlement date, the company pays refunds and handling
costs, extinguishes the remaining return rights, and receives a recovered lot.
Cash, income, liabilities and inventory change together.

Begin with the complete preceding record. The concrete native-state profile
specifies the outer record and every nested component, lot, balance, tax premise,
cash flow and prior-period record. It rejects a missing field, an unexpected
field, a noncanonical rational or a non-Boolean disposition flag before projecting
selected values into a requirement. Fields that can retain an explicit unknown
are named by the schema; unknown tax findings need not prevent an earlier
non-tax event. Structural validity establishes that the record has the specified
meaning and shape. Balance, allocation, financial domains and source authority
remain separate questions.

\subsection{Derive the amounts before changing the book}

Suppose the company accepts two returned units. Each unit entitles the customer
to a refund of USD~20. Its original unit cost is USD~12, deterioration is USD~2
per unit, and handling costs USD~4 per unit. The supplied net realisable value
of the recovered lot is USD~20. The selected method treats handling as a selling
expense and requires an explicit finding that all return rights are closed.
These are the assumptions of this example, including the specified treatment
of the recovered stock; a different accounting method needs its own definition.

Let $n$ be the returned quantity, $p$ the per-unit refund, $c$ the original unit
cost, $d$ the per-unit deterioration, $h$ the handling cost per unit and $N$ the
lot's net realisable value. All monetary inputs refer to the same currency and
selected measurement context. The refund $Q$, handling payment $H$ and recovered
carrying amount $V$ are
\begin{equation}
 Q=np=40,\qquad H=nh=8,\qquad
 V=\min\{n(c-d),N\}=\min\{2(12-2),20\}=20.
\label{eq:typed-return-amounts}
\end{equation}
Units multiplied by currency per unit produce currency. Subtracting $c-d$
requires matching units, scale and context; taking the minimum requires two
amounts with the same full type. If $N$ were USD~5 under a separately supplied
scenario, the carrying amount would be USD~5. That changed scenario would no
longer be the exact retained source event.

The preceding refund liability has a positive magnitude $F=40$, and the
preceding recovery asset is $R=16$. These are balances carried from the preceding
state, not estimates recalculated at settlement. With debits positive, the
selected journal is described by
\begin{equation}
\begin{aligned}
 \Delta\mathrm{RefundLiability}&=F,&
 \Delta\mathrm{RecoveryAsset}&=-R,&
 \Delta\mathrm{Inventory}&=V,\\
 \Delta\mathrm{Revenue}&=Q-F,&
 \Delta\mathrm{CostOfSales}&=R-V,&
 \Delta\mathrm{SellingExpense}&=H,\\
 \Delta\mathrm{Cash}&=-Q-H.
\end{aligned}
\label{eq:typed-return-journal}
\end{equation}
Thus the liability receives a debit of~40, the recovery asset a credit of~16,
and inventory a debit of~20. Revenue has no net adjustment; cost of sales is
credited by~4, selling expense debited by~8, and cash credited by~48. Adding the
signed entries cancels $F$, $R$, $V$, $Q$ and $H$ exactly. The cash-flow record
retains the refund outflow of~40 and handling outflow of~8 separately as operating
flows. Netting them into one number would lose their distinct origins.

\code{MethodMonetaryEffects/1} expresses these calculations as a closed program
of named inputs, ordered calculations, signed postings and gross flows. It
reuses the dimensional product rules of the requirement language. Addition,
subtraction, multiplication and binary minimum evaluate exact rational
coefficients; no implicit rounding, scale conversion or currency conversion is
performed. A calculation can reference an original input or an earlier named
calculation. Its declared dependencies identify the original inputs even when
intermediate names are used.

Unknown values remain unknown through this arithmetic. If $N$ is unavailable,
the program retains that reason and withholds a complete monetary proposal.
Independent known calculations, such as the refund, remain available for
diagnosis. Even multiplication by zero does not silently turn an unobserved
input into an observed fact. The profile also checks finite expression, record
and coefficient limits; exceeding a limit produces an explicit rejection.

Balance and cash agreement are necessary checks, but they do not establish the
measurement. Replacing $V$ by $V+1$ in both affected postings still gives a
balanced journal. The inventory and cost-of-sales amounts would nevertheless
be wrong relative to equation~\eqref{eq:typed-return-amounts}. The implementation
therefore compares the derived amounts and complete effects, rather than using
journal balance as a substitute for the method.

\subsection{Stage the whole event or retain the preceding state}

Settlement also changes memory. The component's refund and recovery amounts
become zero and its settled flag becomes true. A new lot records the two units,
carrying amount~20, acquisition date and component identity. In this selected
tax treatment its tax base equals its carrying amount, and the tax adjustment
is $R-F=-24$. Those tax choices are retained conditional premises. They are not
inferred from the mere existence of the accounting journal.

\code{MethodMemoryEffects/1} provides three operations over closed typed records:
replace an existing object member, insert an absent map member, and append an
array item. A value may be a literal, an exact read from the preceding state or
event-derived inputs, or a constructed record or array. Reads always use the
original records; they cannot accidentally observe an earlier staged write.
The complete preceding and resulting states must satisfy the selected schema.

Paths use explicit tokens. A lot identifier supplied through an input is one
map key, including any slash or tilde that the applicable schema permits. It
cannot redirect an update into another branch of the record. Overlapping writes
are rejected, apart from repeated appends to the same array, whose supplied
order is retained. A duplicate lot insertion, missing parent, wrong value type
or failed final schema check discards the staged result. The caller's original
state remains unchanged.

For a subsequent whole-lot sale with proceeds $P$ and carrying amount $V$, the
selected method requires completed performance and a live, undisposed lot. It
posts cash $P$, revenue $-P$, inventory $-V$ and cost of sales $V$. The lot remains
in memory with zero units, carrying amount and tax base, and a true disposition
flag. Its identity cannot be reused for another sale. Events with equal dates
retain supplied order: a sale before its lot is created fails, while a sale
after settlement on that same date can satisfy the identity condition.

The complete transition checks the preceding book, evaluates the event's
requirements, derives monetary effects, adds them to balances, stages memory
and validates the resulting book. Cash below zero is a funding failure under
this profile's absence of an additional financing policy. It is not repaired
by inventing borrowing or by committing only the noncash entries. The preceding
state and all its identity records remain available after a rejected event.

\subsection{Recompute the proposal from its inputs}

The selected \code{TypedReturnTransition/1} adapter binds exact raw state and
event bytes to the two recorded programs. Its guarded requirements retain
original input paths and raw identities. Dates must lie inside the reporting
period and cannot precede either the state date or the tax close. Refund and
handling cash facts must equal the quantities derived by the method. A false
rights or performance finding is outside the selected method; an unavailable
required finding remains unresolved.

The adapter carries preceding state amounts into the declared event measurement
context. This is an explicit interpretation of the selected transition: the
original state date and source path remain recorded. Equality of type contexts
does not itself establish that information was available at the event date.
The concrete profile is restricted to the retained USD legal book and the two
selected operations. Other currencies and tax or surrounding operations require
their own admitted method routes.

A proposal includes the raw and decoded input identities, selected program and
schema identities, prerequisite findings, journal, gross flows and complete
after-state. Applying it recomputes every field from the original inputs before
returning the result. A modified journal, cash-flow origin, lot flag, input
identity or extra proposal field fails that comparison. A matching hash supplied
by a caller cannot replace the computation. Applying the same event to its
resulting state also fails the lifetime-event uniqueness condition.

The reference adapter reuses the frozen native invariant before and after the
transition. Agreement with the native implementation therefore provides an
independent check of the typed effect derivation, but not an independent proof
of that shared invariant. The separate state-requirement examples provide
additional, explicitly scoped evidence. Where several requirements fail,
the common assessment retains all findings and may summarize them differently
from the native function's first-error order.

\subsection{Correspondence with a source is a separate check}

A source pointer with the right syntax does not show that an event was reported
or that its preceding book follows from the source. The retained example has an
additional finite binder. It accepts the exact original synthetic packet,
checks the schema-prefixed identities of its records, and reconstructs the
prescribed event projection. Legacy exact numbers such as \code{2} become
canonical rationals such as \code{2/1}; both the original packet identity and
the derived record identity remain available.

To establish the preceding book for a selected event, the binder initializes the
opening and recomputes the prior events in the packet's supplied order. It uses
the pinned native and surrounding financial methods, without adopting saved
after-states or expected journals. It then requires exact agreement between that
recomputed state, the prescribed event and the supplied transition inputs. A
balanced book with an extra unit of cash offset by capital is consequently
rejected: balance alone does not establish correspondence with this history.

The projection includes authored interface identifiers and conditional tax
premises. Its zero opening tax bases follow the retained taxable-history
assumption; they are not observations extracted from a company's filing. The
packet also contains nonfinancial build timing, handled by a separate bounded
parser for these exact bytes. The strict financial-record decoder continues to
reject floating monetary inputs and ambiguous encodings.

This chain supplies a reproducible transition for the selected synthetic case:
complete state structure, explicit prerequisites, typed amounts, atomic memory,
full proposal replay and exact source correspondence. It leaves accounting and
source-authority admission false, and its in-memory application makes no durable
storage or production-concurrency claim. Reporting, annual carry, alternative
histories and real-company feasibility retain their distinct obligations. The
checked implementation packages and their exact limitations are preserved in
the accompanying transition and source-binding evidence records.
